% !TeX root = ../../Book1.tex
% 纲要标题：## 第18章　嵌入的独立性、迹与范数
% 纲要位置：计划纲要.md，第 614 行。

\chapter{嵌入的独立性、迹与范数}
\label{b1:ch:18}
\BookChapterNumber{b1:ch:18}

\begin{chapterabstract}
本章用嵌入和线性代数考察有限扩张。Artin 独立性从最短线性关系出发，证明不同乘法特征及域嵌入作为函数线性无关。乘法算子的迹和行列式随后给出域迹与域范数；共轭公式、塔式公式和迹配对把可分性转化为可计算的条件。最后讨论有限域中的像与核，并就地建立整元素的基本封闭性，为数域算术中的迹、范数和判别式作准备。
\end{chapterabstract}

\begin{learninggoals}
\begin{itemize}
\item 掌握最短关系法，区别乘法特征的线性独立性与表示特征标的正交性。
\item 从乘法矩阵与最小多项式计算迹、范数，正确处理不可分重数。
\item 证明迹与范数的传递性，用迹配对和判别式检测有限扩张的可分性。
\item 计算有限域上的迹与范数，说明整性如何保证数域中的相应值为整数。
\end{itemize}
\end{learninggoals}

本章使用域扩张的次数、嵌入计数及可分次数的塔式公式，并假定读者熟悉有限维线性代数中的矩阵迹、行列式与换基。Artin 独立性不依赖 Galois 对应，反而将为下一章的不动域定理提供主要工具。整性所需的有限生成模论证在末节完整给出，不把第二、三卷尚未建立的结构定理当作前提。

\ProseParagraphBreak
迹、范数和乘法算子的关系，可结合 Roman 的《Field Theory》\cite[Section 8.1]{Roman2006FieldTheory}阅读；整性可继续查阅 Atiyah 与 Macdonald 的《Introduction to Commutative Algebra》\cite[Chapter 5]{AtiyahMacdonald1969CommutativeAlgebra}。本章还区分乘法特征与表示的特征标；要学习后一种概念，可读 Serre 的《Linear Representations of Finite Groups》\cite[Part I, Chapters 1--2, pp. 3--24]{Serre1977FiniteRepresentations}。

% !TeX root = ../../Book1.tex
% 纲要标题：### 18.1 Artin 独立性
% 纲要位置：计划纲要.md，第 616 行。

\section{Artin 独立性}
\label{b1:sec:1801}

% 纲要内容（仅作写作计划，尚非教材正文）：
%
% * 乘法特征的线性无关性
% * 域嵌入的线性无关性
% * 与有限群表示特征标正交性的区别
% * 先以最短非平凡线性关系证明任意群上不同乘法特征的有限线性独立性，再得到不同域嵌入的线性独立性；本节不依赖 Galois 对应，且区分函数空间中的线性关系与复特征标内积
%

两个嵌入即使有相同的定义域和值域，也可能作为函数完全不同。我们需要的结论比“两个映射不相等”强：不同嵌入之间不存在系数恒定的非平凡线性关系。证明只依赖乘法，因此先在任意群上建立它。

\subsection{乘法特征的线性无关性}
\label{b1:subsec:1801:01}

设 \(G\) 为群，\(\Omega\) 为域。保持乘法的映射 \(\chi:G\rightarrow\Omega^\times\) 称为一个 \(\Omega\) 值的\term[chengfa-tezheng]{乘法特征}[multiplicative character]。把这些映射视为全体函数组成的 \(\Omega\) 向量空间中的元素，才能谈论线性无关；关系 \(\sum_i c_i\chi_i=0\) 要求同一组常数 \(c_i\) 对每个 \(g\in G\) 都使函数值之和为零。

\ProseParagraphBreak
以下结果称为\term[Artin-dulixing]{Artin 独立性定理}[Artin independence theorem]。
\begin{theorem}[Artin]\label{b1:thm:artin-independence}
设 \(G\) 为群，\(\Omega\) 为域。从 \(G\) 到 \(\Omega^\times\) 的任意一族互异群同态，在全体函数 \(G\rightarrow\Omega\) 组成的 \(\Omega\) 向量空间中线性无关。对无限族，这表示每个有限子族都线性无关。
\end{theorem}
\begin{proof}
若有非平凡关系，从全部这种关系中取非零项数最少者，重新编号写为
\[
c_1\chi_1(g)+\cdots+c_r\chi_r(g)=0\qquad(g\in G),
\]
其中每个 \(c_i\ne0\)。因为 \(\chi_1(g)\ne0\)，不可能 \(r=1\)。由 \(\chi_1\ne\chi_r\)，存在 \(h\in G\) 使 \(\chi_1(h)\ne\chi_r(h)\)。将关系中的 \(g\) 换成 \(hg\)，再减去原关系的 \(\chi_r(h)\) 倍，得到
\[
\sum_{i=1}^{r-1}c_i\Bigl(\chi_i(h)-\chi_r(h)\Bigr)\chi_i(g)=0
\qquad(g\in G).
\]
第一项系数非零，而项数至多为 \(r-1\)，与最小性矛盾。全部操作只涉及有限个项，所以同样排除了无限族中的有限非平凡关系。
\end{proof}
群 \(G\) 不必有限，也不必交换；域的特征没有限制。证明中没有对群元素求平均，也没有除以群阶。

\subsection{域嵌入的线性无关性}
\label{b1:subsec:1801:02}

\begin{corollary}\label{b1:cor:embedding-independence}
设 \(L,\Omega\) 为域。若 \(\sigma_1,\ldots,\sigma_r:L\rightarrow\Omega\) 为互异的域嵌入，则它们在全体函数 \(L\rightarrow\Omega\) 组成的 \(\Omega\) 向量空间中线性无关。
\end{corollary}
\begin{proof}
将嵌入限制到 \(L^\times\)，得到乘法特征。不同域嵌入必在某个非零元素上不同，因为它们都把零映为零。因此这些限制互异，\cref{b1:thm:artin-independence}排除了其上的线性关系，也就排除了原映射的线性关系。
\end{proof}
例如复数上的恒等映射与共轭映射线性无关：若 \(az+b\bar z=0\) 对所有 \(z\in\mathbb C\) 成立，分别代入 \(z=1,i\)，就有 \(a+b=0,a-b=0\)，从而 \(a=b=0\)。一般域上未必有这样显眼的一对检验点，独立性定理保证仍能找到足够多的点。

\subsection{乘法特征与表示特征标}
\label{b1:subsec:1801:03}

本节的乘法特征给出一维线性表示。一般有限群表示的特征标，是给每个群元素的表示矩阵取迹后得到的函数；维数大于一时，这样的迹函数通常不保持乘法。第五卷将研究复数域上不可约表示特征标的正交性，届时需要一个在整个有限群上求和的内积。

\ProseParagraphBreak
这里没有引入那种内积，也没有把“内积为零”和“函数线性无关”混为同一陈述。即使 \(G\) 有限而 \(\Omega\) 的特征整除 \(|G|\)，以 \(|G|\) 作分母的平均运算无意义，本节结论仍然成立。例如特征 \(p\) 的域中，循环 \(p\) 群到乘法群的特征只能是平凡特征，因为 \(x^p-1=(x-1)^p\)；结论并不凭空制造 \(p\) 个不同特征。

\subsection{最短线性关系的证明方法}
\label{b1:subsec:1801:04}

最短关系法先把一个未知的关系压缩到不能再短，再借助结构运算消去一项。关键在于消去某项的同时，必须明确指出另一个系数没有消失；否则新关系可能只是恒等的零关系。下面把函数的独立性转成后文可以使用的矩阵。

\begin{lemma}\label{b1:lem:independent-functions-evaluation}
设 \(X\) 为集合，\(\Omega\) 为域，函数 \(f_1,\ldots,f_r:X\rightarrow\Omega\) 在全体函数 \(X\rightarrow\Omega\) 组成的 \(\Omega\) 向量空间中线性无关。则存在 \(x_1,\ldots,x_r\in X\)，使矩阵 \(\Bigl(f_i(x_j)\Bigr)_{i,j}\) 可逆。
\end{lemma}
\begin{proof}
考虑列向量 \(v(x)=\Bigl(f_1(x),\ldots,f_r(x)\Bigr)^{\mathsf T}\) 在 \(\Omega^r\) 中的张成空间。若其维数小于 \(r\)，对该空间的一组基解齐次线性方程，存在非零行向量 \((c_1,\ldots,c_r)\) 在整个空间上取零值。于是 \(\sum_i c_if_i(x)=0\) 对所有 \(x\) 成立，违背独立性。因此这些列向量张成 \(\Omega^r\)，可以从中选出 \(r\) 个组成一组基，所对应的求值矩阵可逆。
\end{proof}
这里的系数属于值域 \(\Omega\)，不必属于嵌入所固定的较小基域。这一点将在不动域定理的次数估计中起作用。

% !TeX root = ../../Book1.tex
% 纲要标题：### 18.2 迹与范数
% 纲要位置：计划纲要.md，第 623 行。

\section{迹与范数}
\label{b1:sec:1802}

% TODO: 按以下纲要要点撰写本节。

% 纲要内容（仅作写作计划，尚非教材正文）：
%
% * 乘法线性算子的迹与行列式
% * 有限扩张中的域迹与域范数
% * 可分情形的共轭元公式
%

\subsection{乘法线性算子的迹与行列式}
\label{b1:subsec:1802:01}

待撰写。

\subsection{有限扩张中的域迹与域范数}
\label{b1:subsec:1802:02}

待撰写。

\subsection{可分情形的共轭元公式}
\label{b1:subsec:1802:03}

待撰写。

% !TeX root = ../../Book1.tex
% 纲要标题：### 18.3 传递性与判别式
% 纲要位置：计划纲要.md，第 629 行。

\section{传递性与判别式}
\label{b1:sec:1803}

% 纲要内容（仅作写作计划，尚非教材正文）：
%
% * 迹、范数的塔式公式
% * 迹配对及可分性的判别
% * 基的判别式及换基公式
%

在一座域扩张塔中，先把乘法压缩到中间域，再压缩到基域，所得迹与范数应当与一步计算相同。迹还可以比较两个元素：对乘积取迹便得到一个双线性配对。这个配对是否退化，恰好检测扩张的可分性。

\subsection{迹、范数的塔式公式}
\label{b1:subsec:1803:01}

\begin{theorem}\label{b1:thm:trace-norm-tower}
若 \(K\subseteq E\subseteq L\) 为有限扩张塔，则对 \(a\in L\) 有
\begin{align*}
\operatorname{Tr}_{L/K}(a)&=\operatorname{Tr}_{E/K}\Bigl(\operatorname{Tr}_{L/E}(a)\Bigr),\\
\operatorname{N}_{L/K}(a)&=\operatorname{N}_{E/K}\Bigl(\operatorname{N}_{L/E}(a)\Bigr).
\end{align*}
\end{theorem}
\begin{proof}
取 \(E/K\) 的基 \(u_1,\ldots,u_m\) 与 \(L/E\) 的基 \(v_1,\ldots,v_r\)，由\cref{b1:thm:tower-law}的基构造，全部 \(u_i v_j\) 为 \(L/K\) 的基。设乘以 \(a\) 的 \(E\) 矩阵为 \(A=(a_{ij})\)。把每个系数 \(a_{ij}\) 换成它作用在 \(E\) 上的 \(K\) 乘法矩阵，便得到乘以 \(a\) 的 \(K\) 矩阵。因此总迹为
\[
\sum_j\operatorname{Tr}_{E/K}(a_{jj})
=\operatorname{Tr}_{E/K}\Bigl(\sum_j a_{jj}\Bigr),
\]
即迹的塔式公式。

对范数，需要说明这个分块矩阵的行列式为 \(\operatorname{N}_{E/K}(\det_E A)\)。若 \(A\) 奇异，非零的 \(E\) 核向量也是底层 \(K\) 线性映射的核向量，两边皆为零。若 \(A\) 可逆，可用三类初等行变换化为单位矩阵。加上另一行的 \(c\) 倍，在 \(K\) 分块矩阵中是单位对角的块初等变换，两种行列式因子均为一。把一行乘以 \(c\ne0\)，块行列式乘以 \(\operatorname{N}_{E/K}(c)\)，与 \(\det_E A\) 乘以 \(c\) 后再取范数相同。交换两行，在 \(K\) 上交换两个各有 \(m\) 个坐标的块，符号为 \((-1)^{m^2}=(-1)^m=\operatorname{N}_{E/K}(-1)\)。于是两种行列式在每步受到相同的因子改变，单位矩阵处又均为一，故所需等式成立。取 \(A=m_a\) 的 \(E\) 矩阵便得到范数公式。这个论证不要求可分性。
\end{proof}

\subsection{迹配对及可分性的判别}
\label{b1:subsec:1803:02}

\begin{definition}\label{b1:def:trace-pairing}
设 \(L/K\) 为有限域扩张。双线性映射
\[
L\times L\longrightarrow K,\qquad (x,y)\longmapsto\operatorname{Tr}_{L/K}(xy)
\]
称为 \(L/K\) 的\term[ji-peidui]{迹配对}[trace pairing]。如果对每个非零 \(x\in L\) 都存在 \(y\in L\) 使 \(\operatorname{Tr}_{L/K}(xy)\neq0\)，则称此配对非退化。
\end{definition}
\begin{theorem}\label{b1:thm:trace-pairing-separable}
有限扩张 \(L/K\) 可分，当且仅当其迹配对非退化。
\end{theorem}
\begin{proof}
若扩张可分，由\cref{b1:thm:trace-norm-embeddings}，迹函数为全部不同 \(K\) 嵌入的和。\cref{b1:cor:embedding-independence}说明这个和不是零函数，因为其各系数均为非零的一。于是存在 \(z\in L\) 使 \(\operatorname{Tr}_{L/K}(z)\ne0\)。对任意 \(x\ne0\)，取 \(y=x^{-1}z\)，便有 \(\operatorname{Tr}_{L/K}(xy)\ne0\)。

若扩张不可分，\cref{b1:prop:separable-inseparable-decomposition}说明其特征为某个素数 \(p\)，不可分次数 \(e\) 是大于一的 \(p\) 的幂。由\cref{b1:thm:trace-norm-embeddings}，迹函数是某个和的 \(e\) 倍，因而恒为零。所以整个迹配对为零，必退化。
\end{proof}
可分扩张的次数即使被特征整除，迹仍不是零映射；此时只是 \(\operatorname{Tr}_{L/K}(1)=[L:K]=0\)。把“单位元的迹为零”与“所有元素的迹都为零”区别开来十分必要。

\subsection{基的判别式及换基公式}
\label{b1:subsec:1803:03}

设 \(L/K\) 为 \(n\) 次有限域扩张，\(b=(b_1,\ldots,b_n)\) 是 \(L\) 在 \(K\) 上的一组基。把
\[
\operatorname{disc}_{L/K}(b)
=\det\biggl(\Bigl(\operatorname{Tr}_{L/K}(b_i b_j)\Bigr)_{i,j}\biggr)
\]
称为基 \(b\) 关于扩张 \(L/K\) 的\term[ji-panbieshi]{基的判别式}[discriminant of a basis]。
\symidx{\(\operatorname{disc}_{L/K}(b)\)：基 \(b\) 的判别式}
若另一组基满足 \(c_j=\sum_i b_i p_{ij}\)，记换基矩阵为 \(P\)。双线性给出迹矩阵关系 \(T_c=P^{\mathsf T}T_bP\)，所以
\begin{equation}\label{b1:eq:discriminant-change-basis}
\operatorname{disc}_{L/K}(c)=(\det P)^2\operatorname{disc}_{L/K}(b).
\end{equation}
由\cref{b1:thm:trace-pairing-separable}，扩张可分当且仅当任意一组基的判别式非零。

\ProseParagraphBreak
在可分情形，令 \(M=\Bigl(\sigma_i(b_j)\Bigr)_{i,j}\) 为嵌入求值矩阵。共轭公式给出 \(T_b=M^{\mathsf T}M\)，故判别式为 \((\det M)^2\)。若 \(L=K(a)\)，基取 \(1,a,\ldots,a^{n-1}\)，各共轭根为 \(a_1,\ldots,a_n\)，Vandermonde 行列式得到
\begin{equation}\label{b1:eq:power-basis-discriminant}
\operatorname{disc}_{L/K}(1,a,\ldots,a^{n-1})
=\prod_{i<j}(a_j-a_i)^2.
\end{equation}
一般地，设 \(K\) 为域，\(f\in K[x]\) 为次数 \(m\) 的首一多项式，并把它在一个分裂域中的全部根按重数列为 \(a_1,\ldots,a_m\)。把 \(\prod_{i<j}(a_j-a_i)^2\) 称为 \(f\) 的\term[duoxiangshi-panbieshi]{判别式}[discriminant of a polynomial]。它与根的排列无关，在根重合时为零；当 \(f\) 是上文中元素 \(a\) 的首一最小多项式时，这一定义就是\eqref{b1:eq:power-basis-discriminant}右端。对 \(\mathbb Q(\sqrt d)\) 的基 \(1,\sqrt d\)，迹矩阵为 \(\operatorname{diag}(2,2d)\)，判别式为 \(4d\)。基的判别式依赖所选基，不能直接把任意一组有理基的判别式称为数域整数环的判别式。

% !TeX root = ../../Book1.tex
% 纲要标题：### 18.4 迹与范数的算术应用
% 纲要位置：计划纲要.md，第 635 行。

\section{迹与范数的算术应用}
\label{b1:sec:1804}

% 纲要内容（仅作写作计划，尚非教材正文）：
%
% * 有限域中的迹和范数
% * 整元素的迹、范数与整性
% * 数域整数环和分歧理论的后续方向
% * 所需整元素的首一方程定义及整性的基本封闭性在本节就地准备；通过共轭元或有限模方法证明迹与范数的整性，不以前置第三卷整扩张理论为条件
%
% ---
%

迹和范数把扩张域中的计算带回基域，但它们的像并非在所有场合都一样。有限域中可以精确求出像与核；数域中则要同时考虑元素是否满足整数系数的首一方程。本节就地建立所需的整性事实，不预先使用第三卷的整扩张理论。

\subsection{有限域中的迹和范数}
\label{b1:subsec:1804:01}

\begin{proposition}\label{b1:prop:finite-field-trace-norm}
对 \(L=\mathbb F_{q^n}\)、\(K=\mathbb F_q\)，有
\[
\operatorname{Tr}_{L/K}(a)=a+a^q+\cdots+a^{q^{n-1}},\qquad
\operatorname{N}_{L/K}(a)=a^{1+q+\cdots+q^{n-1}}.
\]
迹作为加法群同态满射到 \(K\)，其核有 \(q^{n-1}\) 个元素；范数把 \(L^\times\) 满射到 \(K^\times\)，其核有 \((q^n-1)/(q-1)\) 个元素。
\end{proposition}
\begin{proof}
\cref{b1:thm:finite-field-automorphisms}列出了全部嵌入 \(a\mapsto a^{q^i}\)，而有限域扩张可分，故\cref{b1:thm:trace-norm-embeddings}给出两公式。\cref{b1:thm:trace-pairing-separable}保证迹非零；从 \(n\) 维 \(K\) 空间到一维 \(K\) 空间的非零线性映射必满射，核维数为 \(n-1\)，从而有 \(q^{n-1}\) 个元素。取 \(L^\times\) 的生成元 \(g\)，令 \(h=(q^n-1)/(q-1)\)，则 \(g^h\) 的阶为 \(q-1\)，所以范数的像包含 \(K^\times\) 的一个生成元，必为全部。循环群中幂映射的核有 \(h\) 个元素，得到最后一项。
\end{proof}
例如在 \(\mathbb F_4=\mathbb F_2(u)\)、\(u^2+u+1=0\) 中，\(\operatorname{Tr}(1)=0\)，但 \(\operatorname{Tr}(u)=u+u^2=1\)；这正展示了迹非零并不要求单位元的迹非零。

\subsection{整元素的迹与范数}
\label{b1:subsec:1804:02}

\begin{definition}\label{b1:def:integral-element-field}
设 \(A\) 是域 \(L\) 的子环。若 \(a\in L\) 满足某个系数在 \(A\) 中的首一方程
\[
a^d+c_{d-1}a^{d-1}+\cdots+c_0=0,
\]
则称 \(a\) 为 \(A\) 上的\term[zheng-yuansu]{整元素}[integral element]，也说它在 \(A\) 上整。在 \(\mathbb Z\) 上整的代数数称为\term[daishu-zhengshu]{代数整数}[algebraic integer]。
\end{definition}
有限扩张 \(K/\mathbb Q\) 称为\term[shuyu]{数域}[number field]。为将整性用于迹与范数，先给出求和、求积时所需的封闭性。
\begin{lemma}\label{b1:lem:integral-closure-elementary}
在同一个域中，有限多个在 \(A\) 上整的元素生成的子环 \(A[a_1,\ldots,a_s]\) 由有限多个元素的 \(A\) 系数线性组合组成，也称有限生成 \(A\) 模；其中每个元素都在 \(A\) 上整。特别地，整元素的和、差与积仍整。
\end{lemma}
\begin{proof}
设 \(a_i\) 满足次数为 \(d_i\) 的首一方程。反复用这些方程约去高次幂，\(A[a_1,\ldots,a_s]\) 由有限多个单项式 \(a_1^{e_1}\cdots a_s^{e_s}\)、\(0\leq e_i<d_i\)，作为 \(A\) 模生成。记这些生成元为 \(v_1,\ldots,v_m\)，其中取 \(v_1=1\)。

若 \(b\) 在这个代数中，将每个 \(bv_i\) 写成生成元的线性组合，得到 \(bv_i=\sum_j c_{ij}v_j\)，其中 \(c_{ij}\in A\)。令 \(C=(c_{ij})\)、\(v=(v_i)^{\mathsf T}\)，则 \((bI-C)v=0\)。由行列式的余子式展开，伴随矩阵满足 \(\operatorname{adj}(bI-C)(bI-C)=\det(bI-C)I\)；乘到上式得到 \(\det(bI-C)v_i=0\) 对每个 \(i\) 成立。取 \(v_1=1\)，便得 \(\det(bI-C)=0\)。多项式 \(\det(tI-C)\in A[t]\) 首一，故 \(b\) 在 \(A\) 上整。此处生成元不必线性无关，因而没有假设这个模是自由模。
\end{proof}
这段有限生成模论证只使用余子式恒等式；行列式技巧的一般模论形式将在第三卷系统讨论。

\begin{proposition}\label{b1:prop:integral-trace-norm}
设 \(L/K\) 为有限扩张，\(A\subseteq K\) 为子环。若 \(a\in L\) 在 \(A\) 上整，则 \(\operatorname{Tr}_{L/K}(a)\) 和 \(\operatorname{N}_{L/K}(a)\) 也是 \(K\) 中在 \(A\) 上整的元素。特别地，数域中的代数整数对 \(\mathbb Q\) 的迹和范数都是整数。
\end{proposition}
\begin{proof}
取共同代数闭包。\(a\) 满足的首一 \(A\) 方程在每个保持 \(K\) 的嵌入下不变，所以每个共轭像也在 \(A\) 上整。由\cref{b1:thm:trace-norm-embeddings}，迹和范数是这些共轭像按不可分重数形成的有限和、积。\cref{b1:lem:integral-closure-elementary}保证它们整，而定义保证它们属于 \(K\)。

若 \(K=\mathbb Q,A=\mathbb Z\)，还须说明有理数中整元素只有整数。写既约分数 \(a/b\)、\(b>0\)，把它满足的首一 \(d\) 次整数方程乘以 \(b^d\)，得到 \(b\mid a^d\)。由 \(\gcd(a,b)=1\)，只能 \(b=1\)。所以迹与范数为整数。
\end{proof}

\subsection{数域整数环与分歧理论}
\label{b1:subsec:1804:03}

数域 \(K\) 中全部代数整数组成的集合记为 \(\mathcal O_K\)，称为其\term[zhengshu-huan]{整数环}[ring of integers]。
\symidx{\(\mathcal O_K\)：数域 \(K\) 的整数环}
\cref{b1:lem:integral-closure-elementary}说明它对加、减、乘封闭并含一，所以确实是环。例如 \(\sqrt d\) 是代数整数，因为满足 \(t^2-d=0\)。当平方自由整数 \(d\equiv1\pmod4\) 时，\(\omega=(1+\sqrt d)/2\) 也满足
\[
t^2-t+\frac{1-d}{4}=0,
\]
故为代数整数；这已说明 \(\mathbb Z[\sqrt d]\) 未必是整个整数环。

\ProseParagraphBreak
迹配对把数域的线性结构与整数环的算术联系起来。例如 \(d=5\) 时，基 \(1,\omega\) 的判别式为五，而基 \(1,\sqrt5\) 的判别式为二十；两者符合\eqref{b1:eq:discriminant-change-basis}。关于整数环的整基、素理想分解，以及判别式和分歧之间的关系，本节只指出研究方向，尚未建立相应的整数环结构定理。特别不能把某组基的判别式中的每个素因子都直接宣布为数域中的分歧素数。

\subsection{整性的定义与封闭性}
\label{b1:subsec:1804:04}

整性的定义不要求系数环为域。事实上，若系数环已经是域 \(K\)，首一方程与代数方程可以通过除以首项系数互相转换，“在 \(K\) 上整”就等于“在 \(K\) 上代数”。对于一般环 \(A\)，首项系数不能随意取逆，整性因此比代数性强。

\begin{proposition}\label{b1:prop:integrality-transitive-elementary}
设域内有子环 \(A\subseteq B\)，且 \(B\) 中每个元素都在 \(A\) 上整。如果 \(c\) 在 \(B\) 上整，则 \(c\) 在 \(A\) 上整。
\end{proposition}
\begin{proof}
取 \(c\) 满足的一条首一方程，其有限多个系数为 \(b_1,\ldots,b_s\)。由\cref{b1:lem:integral-closure-elementary}，\(B_0=A[b_1,\ldots,b_s]\) 有有限个 \(A\) 模生成元。若该方程次数为 \(d\)，则 \(B_0[c]\) 由 \(1,c,\ldots,c^{d-1}\) 作为 \(B_0\) 模生成；把两组生成元相乘，得到有限个 \(A\) 模生成元，并可令其中一个为一。对乘以 \(c\) 再应用\cref{b1:lem:integral-closure-elementary}证明中的伴随矩阵论证，便得到 \(c\) 的首一 \(A\) 方程。
\end{proof}
整元素的逆未必整，例如二是代数整数而 \(1/2\) 不是；迹、范数均为整数也不能在任意次数下反推元素整，因为它们只控制特征多项式的两个系数。对于二次扩张，这两个系数已经是全部非首项系数，所以结论恰好可以反推；章末习题将比较这两种情形。


\begin{chaptersummary}
不同域嵌入之间不存在常系数的非平凡线性关系，这使可分扩张的迹成为非零映射。迹与范数首先是乘法算子的线性代数量；共轭和积只是它们在代数闭包中的表达，不可分时必须计入统一重数。塔式公式适用于所有有限扩张，迹配对非退化则恰等价于可分。判别式记录该配对在一组基下的行列式，并随换基乘以一个平方。整性封闭性进一步把这些运算限制到算术对象中，整数环的结构及分歧理论将在后续内容中讨论。
\end{chaptersummary}

\BookExercises

\begin{exercise}\label{b1:ex:18-characters-integers}
设 \(a_1,\ldots,a_r\) 为域 \(\Omega\) 中不同的非零元素。对整数 \(n\)，令 \(\chi_i(n)=a_i^n\)。证明这些函数线性无关，并说明只检验 \(n=0,\ldots,r-1\) 为何已经足够。
\end{exercise}
\begin{solution}
它们是加法群 \(\mathbb Z\) 到 \(\Omega^\times\) 的不同乘法特征，可直接用\cref{b1:thm:artin-independence}。若只取指定的 \(r\) 个整数，求值矩阵为 Vandermonde 矩阵，其行列式为 \(\prod_{i<j}(a_j-a_i)\ne0\)。因此在这些点成立的 \(r\) 个齐次方程已经迫使全部系数为零。
\end{solution}

\begin{exercise}\label{b1:ex:18-evaluation-matrix}
在 \(L=\mathbb Q(\sqrt2,\sqrt3)\) 中，以 \(1,\sqrt2,\sqrt3,\sqrt6\) 为基，将独立改变两个平方根符号的四个嵌入作为矩阵的行。求求值矩阵的行列式的平方，并据此求该基的判别式。
\end{exercise}
\begin{solution}
\cref{b1:prop:biquadratic-example}给出扩张次数为四及题中所列的基。取行的符号依次为 \((+,+)\)、\((-,+)\)、\((+,-)\)、\((-,-)\)，矩阵可写成
\[
\begin{pmatrix}1&1&1&1\\1&-1&1&-1\\1&1&-1&-1\\1&-1&-1&1\end{pmatrix}
\operatorname{diag}(1,\sqrt2,\sqrt3,\sqrt6).
\]
左矩阵与其转置的乘积为 \(4I_4\)，故行列式平方为 \(4^4=256\)；右矩阵的行列式为六。总行列式平方为 \(256\cdot36=9216\)。扩张可分，迹矩阵为求值矩阵转置乘自身，所以基的判别式也为 \(9216\)。
\end{solution}

\begin{exercise}\label{b1:ex:18-cubic-trace-norm}
令 \(a=\sqrt[3]2\)，\(L=\mathbb Q(a)\)。求 \(a\)、\(a^2\)、\(1+a\) 对 \(\mathbb Q\) 的迹与范数。
\end{exercise}
\begin{solution}
\(a\) 的最小多项式为 \(t^3-2\)，由\eqref{b1:eq:trace-norm-minimal-polynomial}得到迹零、范数二。元素 \(a^2\) 仍生成 \(L\)，因为 \(a=(a^2)^2/2\)；其最小多项式为 \(t^3-4\)，故迹零、范数四。\(1+a\) 也生成 \(L\)，最小多项式为 \((t-1)^3-2=t^3-3t^2+3t-3\)，故迹三、范数三。
\end{solution}

\begin{exercise}\label{b1:ex:18-ambient-extension}
比较 \(\sqrt2\) 在 \(\mathbb Q(\sqrt2)/\mathbb Q\) 与 \(\mathbb Q(\sqrt2,\sqrt3)/\mathbb Q\) 中的迹、范数和乘法特征多项式。
\end{exercise}
\begin{solution}
前一个扩张的最小多项式与特征多项式均为 \(t^2-2\)，迹为零，范数为负二。后一个扩张对 \(\mathbb Q(\sqrt2)\) 的次数为二，因此\eqref{b1:eq:trace-norm-minimal-polynomial}给出特征多项式 \((t^2-2)^2\)，迹仍为零，范数变为四。元素本身和最小多项式都未变，但所处的乘法空间维数不同。
\end{solution}

\begin{exercise}\label{b1:ex:18-not-ring-maps}
举例说明域迹一般不保持乘法，域范数一般不保持加法。说明它们各自确实保持哪些运算。
\end{exercise}
\begin{solution}
在 \(\mathbb C/\mathbb R\) 中，\(\operatorname{Tr}(i)=0\)，而 \(\operatorname{Tr}(i^2)=\operatorname{Tr}(-1)=-2\)，故迹不保持乘法；\(\operatorname{N}(1+1)=4\)，但 \(\operatorname{N}(1)+\operatorname{N}(1)=2\)，故范数不保持加法。\cref{b1:prop:trace-norm-basic}给出的正确结论是：迹为基域线性映射，范数为乘法映射，并在非零元素的乘法群之间给出群同态。
\end{solution}

\begin{exercise}\label{b1:ex:18-tower-calculation}
在 \(L=\mathbb Q(\sqrt2,\sqrt3)\) 中，沿中间域 \(E=\mathbb Q(\sqrt2)\) 计算 \(a=1+\sqrt2+\sqrt3\) 对 \(\mathbb Q\) 的迹与范数。
\end{exercise}
\begin{solution}
在 \(L/E\) 中共轭把 \(\sqrt3\) 变号，所以
\[
\operatorname{Tr}_{L/E}(a)=2+2\sqrt2,\qquad
\operatorname{N}_{L/E}(a)=(1+\sqrt2)^2-3=2\sqrt2.
\]
再对 \(E/\mathbb Q\) 计算，前者迹为四，后者范数为 \((2\sqrt2)(-2\sqrt2)=-8\)。由\cref{b1:thm:trace-norm-tower}，这就是 \(a\) 在整个四次扩张中的迹与范数。
\end{solution}

\begin{exercise}\label{b1:ex:18-purely-inseparable}
设 \(u\) 为超越元，\(L=\mathbb F_p(u)\)、\(K=\mathbb F_p(u^p)\)。证明对每个 \(a\in L\)，有 \(\operatorname{Tr}_{L/K}(a)=0\)、\(\operatorname{N}_{L/K}(a)=a^p\)。
\end{exercise}
\begin{solution}
每个有理函数的 \(p\) 次幂属于 \(K\)，扩张为 \(p\) 次纯不可分扩张。到代数闭包的唯一 \(K\) 嵌入为包含映射，因为 \(u\) 的像必须满足 \(t^p-u^p=(t-u)^p=0\)。\cref{b1:thm:trace-norm-embeddings}中的 \(s=1,e=p\)，故迹为 \(pa=0\)，范数为 \(a^p\)。这也适用于 \(a=0\)。
\end{solution}

\begin{exercise}\label{b1:ex:18-trace-eight}
在 \(\mathbb F_8=\mathbb F_2(\alpha)\)、\(\alpha^3+\alpha+1=0\) 中，列出对 \(\mathbb F_2\) 的迹为零的全部元素，并求每个非零元素的范数。
\end{exercise}
\begin{solution}
\(\operatorname{Tr}(1)=1\)，而 \(\operatorname{Tr}(\alpha)=\alpha+\alpha^2+\alpha^4=0\)，\(\operatorname{Tr}(\alpha^2)=\operatorname{Tr}(\alpha)^2=0\)。因此 \(a+b\alpha+c\alpha^2\) 的迹为 \(a\)，核正是 \(0,\alpha,\alpha^2,\alpha+\alpha^2\)。非零元的范数为其七次幂，均等于一。
\end{solution}

\begin{exercise}\label{b1:ex:18-trace-dual-basis}
设 \(L=\mathbb Q(\sqrt d)\)，\(d\) 为非平方整数。求基 \(1,\sqrt d\) 关于迹配对的对偶基，即求 \(c_1,c_2\) 使 \(\operatorname{Tr}(b_i c_j)=\delta_{ij}\)。
\end{exercise}
\begin{solution}
迹矩阵为 \(\operatorname{diag}(2,2d)\)，所以对偶基为 \(1/2,\sqrt d/(2d)\)。直接相乘：第一、第一的迹为一，两个交叉项均为 \(\sqrt d\) 的有理倍数，迹为零；第二、第二的积为二分之一，迹为一。
\end{solution}

\begin{exercise}\label{b1:ex:18-cubic-discriminant}
对 \(a=\sqrt[3]2\)，直接用迹矩阵计算基 \(1,a,a^2\) 的判别式。
\end{exercise}
\begin{solution}
由 \(a^3=2\) 以及 \(\operatorname{Tr}(a)=\operatorname{Tr}(a^2)=0\)，有 \(\operatorname{Tr}(1)=3\)、\(\operatorname{Tr}(a^3)=6\)、\(\operatorname{Tr}(a^4)=2\operatorname{Tr}(a)=0\)。所以迹矩阵为
\[
\begin{pmatrix}3&0&0\\0&0&6\\0&6&0\end{pmatrix},
\]
行列式为 \(-108\)。非零性符合特征零下的可分性，负号本身不表示配对退化。
\end{solution}

\begin{exercise}\label{b1:ex:18-integral-unit}
设 \(K\) 为数域，\(a\in K\) 为非零代数整数。证明 \(a^{-1}\) 也是代数整数，当且仅当 \(\operatorname{N}_{K/\mathbb Q}(a)=\pm1\)。
\end{exercise}
\begin{solution}
若二者均为代数整数，由\cref{b1:prop:integral-trace-norm}，它们的范数是互为倒数的整数，只能各为正一或负一。反向先说明 \(a\) 的乘法特征多项式系数是整数：其根是 \(a\) 的全部嵌入像，各自在 \(\mathbb Z\) 上整；系数是这些根的初等对称和，依\cref{b1:lem:integral-closure-elementary}仍整，而本来属于 \(\mathbb Q\)，所以都是整数。若范数为 \(\pm1\)，这个首一多项式的常数项为 \(\pm1\)。把其在 \(a\) 处的等式除以 \(a^{[K:\mathbb Q]}\)，再乘以该常数项的符号，便得到 \(a^{-1}\) 满足的首一整数方程，故其整。
\end{solution}

\begin{exercise}\label{b1:ex:18-quadratic-integrality}
设 \(K/\mathbb Q\) 为二次扩张，\(a\in K\)。证明 \(a\) 为代数整数当且仅当它的迹、范数均为整数。以 \(a=(1+\sqrt{13})/2\) 为例求出二者和基 \(1,a\) 的判别式。
\end{exercise}
\begin{solution}
必要性由\cref{b1:prop:integral-trace-norm}。反向，\(a\) 满足其乘法特征多项式 \(t^2-\operatorname{Tr}(a)t+\operatorname{N}(a)\)：由\eqref{b1:eq:trace-norm-minimal-polynomial}，这个多项式是其最小多项式的正整数次幂。若两系数为整数，便为首一整数方程。所给元素满足 \(t^2-t-3=0\)，迹为一，范数为负三。两个共轭之差为 \(\sqrt{13}\)，故基判别式为十三。
\end{solution}

\begin{exercise}\label{b1:ex:18-integrality-counterexample}
设 \(a\) 为 \(t^3-\tfrac12t-1\) 的一个根。证明 \(a\) 对 \(\mathbb Q\) 的迹、范数均为整数，但 \(a\) 不是代数整数。
\end{exercise}
\begin{solution}
整数多项式 \(2t^3-t-2\) 的有理根若存在，只可能为 \(\pm1,\pm2,\pm1/2\)；逐一代入都非零，所以它在 \(\mathbb Q\) 上不可约。因此题给首一式就是最小多项式，迹为零，范数为一。若 \(a\) 整，它的全部共轭也整，其初等对称和由\cref{b1:lem:integral-closure-elementary}仍整；这些和为有理数，必须为整数。然而最小多项式的一次项系数为 \(-1/2\)，矛盾。迹与范数在三次情形不能确定全部系数。
\end{solution}

\begin{exercise}\label{b1:ex:18-integrality-tower}
设 \(b=\sqrt2\)，\(a^2=1+b\)。不用数域整数环的结构定理，证明 \(a\) 和 \(a^{-1}\) 都是代数整数，并写出各自的一条首一整数方程。
\end{exercise}
\begin{solution}
\(b\) 整，而 \(a\) 满足 \(t^2-(1+b)=0\)，由\cref{b1:prop:integrality-transitive-elementary}知 \(a\) 整。具体消去 \(b\) 得 \((a^2-1)^2=2\)，即 \(a^4-2a^2-1=0\)。常数项为负一，\(a\ne0\)。除以 \(a^4\) 后可写为 \((a^{-1})^4+2(a^{-1})^2-1=0\)，所以逆元也整。这里没有用到任何整基或唯一分解结论。
\end{solution}
